\ncdlesson{8}{Janteloven i det moderne samfund}{现代社会中的《詹代法则》}
{DU3 Modul 3（B1）｜对标 NCE 2 后半 \& Modul 3 \textit{Læsning Opgave 3 (Tekstudfyldning)} \& \textit{Samfund og Værdier}}
{\textbf{情态动词的主观推测与传闻用法（Epistemisk modalitet）}：高度肯定推测 \textit{må (have)} · 可能性推测 \textit{kan (ikke)} · 传闻转述 \textit{skal / skulle (efter sigende) være} · 理应推断 \textit{bør / burde}}
{在 Modul 3 谈论社会现象与职场文化时，直接下绝对断言往往不够客观。学会用 \textbf{\textit{må være}}（一定是）、\textbf{\textit{kan være}}（可能是）和 \textbf{\textit{skulle efter sigende være}}（据说是）来表达推断与审慎转述，既能显著提升语体成熟度，也是阅读理解判断作者立场的必考线索！}

% ---------------------------------------------------------------------
\begin{foerduleaeser}
带着下面两个问题阅读 Tekst A 和 Tekst B：
\begin{enumerate}
  \item \textit{Hvilke fordele og ulemper kan Janteloven have på arbejdspladsen ifølge Tekst A?}（根据 Tekst A，《詹代法则》在职场上可能带来哪些利与弊？）
  \item \textit{Hvordan oplevede skribenten i Tekst B sit første projektmøde i et dansk team, og hvad skyldtes den rolige stemning?}（在 Tekst B 中，作者在丹麦团队的第一次项目会议体验如何？那种平静的氛围可能源于什么？）
\end{enumerate}
\end{foerduleaeser}

% ---------------------------------------------------------------------
\ncdsubsec{二、Tekst A \& Tekst B（文化深读与 Modul 3 真题场景）}

\subsection*{Tekst A：\textit{Janteloven i det moderne samfund}（文化深读短文 · 约 240 词）}

\begin{tipbox}{语篇说明}
1933 年作家 Aksel Sandemose 写下著名的《詹代法则》（\textit{Janteloven}），深刻影响了北欧的平等与低调文化。阅读时请重点留意\textbf{\color{NCDNavy}深蓝色标出的主观推测情态动词}（\textit{må, kan, vil, bør}）与\textbf{\color{NCDRed}深红色标出的传闻与否定推测用法}（\textit{skulle efter sigende, kan ikke}）。
\end{tipbox}

I 1933 skrev den dansk-norske forfatter Aksel Sandemose romanen \textit{En flygtning krydser sit spor}, hvor han formulerede ti regler kendt som ”Janteloven”. Den første og mest berømte regel lyder: ”Du skal ikke tro, at du er noget.” Selvom bogen blev skrevet for næsten hundrede år siden, \da{kan man stadig mærke} dens indflydelse i dagens Danmark.

Når udlændinge hører om Janteloven første gang, \da{må det være} svært helt at forstå, hvorfor beskedenhed fylder så meget. For mange danskere \da{må lighed være} en af samfundets vigtigste værdier, og man ønsker ikke, at nogen praler af sin succes eller ser ned på andre. På den ene side \da{kan denne holdning skabe} et trygt arbejdsmiljø, hvor chefen og medarbejderne behandler hinanden ligeværdigt. På den anden side \da{kan Janteloven også være} en bremse, hvis dygtige unge bliver bange for at skille sig ud.

I dag diskuterer man meget, om de gamle regler stadig gælder. Janteloven \dabold{skulle efter sigende være} blevet svagere i de store byer, især blandt unge iværksættere og på sociale medier. Flere sociologer peger på, at kulturen \da{må have ændret sig}, fordi man nu i højere grad opmuntrer hinanden til at være stolte af gode resultater. Alligevel \dabold{kan det ikke passe}, at Janteloven er helt forsvundet. Hvis man praler for meget, \da{vil} reaktionen ofte \da{være} kølig stilhed. Derfor \da{bør} man som ny i Danmark både turde vise sine evner og samtidig huske respekten bag den danske ligeværdighed.

\begin{transbox}
\small
\textbf{\color{NCDTeal}[Tekst A 中文精译]}\\
1933 年，丹麦—挪威作家阿克塞尔·桑德莫斯写下了小说《逃亡者跨越自己的足迹》，在书中提出了被称为“詹代法则”的十条规则。第一条也是最著名的一条是：“不要以为你有什么了不起。”尽管这本书写于将近一百年前，人们在今天的丹麦依然能感受到它的影响。

外国人初次听说《詹代法则》时，一定会觉得很难完全理解为什么谦虚会占据这么大的分量。对许多丹麦人来说，平等必然是社会最重要的价值观之一，大家不希望有人炫耀自己的成功或轻视他人。一方面，这种态度可以营造出安全融洽的工作氛围，让老板和员工彼此平等相待；另一方面，如果优秀的年轻人因此不敢脱颖而出，《詹代法则》也可能成为一种阻力。

如今人们经常讨论这些老规矩是否依然适用。据说《詹代法则》在大城市里已经变弱了，尤其是在年轻创业者群体和社交媒体上。多位社会学家指出，这种文化一定已经发生了变化，因为现在人们越来越鼓励彼此为好成绩感到自豪。尽管如此，说《詹代法则》已经彻底消失是不可能成立的。如果一个人吹嘘过头，周围的反应往往会是一片冷淡的沉默。因此，初来丹麦的人既应该敢于展示自己的能力，同时也要记住这份丹麦式平等背后的相互尊重。
\end{transbox}

\subsection*{Tekst B：\textit{Debatindlæg: Mit første projektmøde i Danmark}（对标 Modul 3 Læsning Opgave 3）}

\textbf{阅读微任务（Læsning Opgave 3 段落逻辑补全）：}\textit{Læs den korte kommentar om dansk arbejdskultur, og vælg den sætning (A--D), der passer bedst i hullet [1].}

\begin{emailbox}{职场文化短评｜Ny i et dansk team}
Da jeg startede i mit danske it-team, troede jeg, at jeg skulle sælge mine idéer meget højlydt for at blive hørt. Jeg havde nemlig læst, at konkurrence er vigtig i internationale firmaer. \underline{\hspace{2.8cm}} \textbf{[1]} På mit allerførste projektmøde lagde jeg mærke til, at min chef talte meget roligt og lyttede lige så meget til praktikanterne som til de erfarne udviklere. Derudover drak alle kaffe sammen i pausen og snakkede om weekendens planer uden at vise sig frem. Det \da{må være} et tydeligt eksempel på, hvordan Janteloven stadig præger samarbejdet. Selvom loven \dabold{skulle efter sigende være} mindre synlig i dag, trives beskedenhed stadig på kontoret.
\end{emailbox}

\noindent\small\textbf{备选句子（A--D）：}\\
\textbf{A.} \textit{Derimod må man aldrig stille spørgsmål til sin chefs beslutninger.}\\
\textbf{B.} \textit{Men sådan fungerer et dansk kontor altså ikke altid.}\\
\textbf{C.} \textit{For vores afdeling skulle efter sigende lukke næste måned.}\\
\textbf{D.} \textit{Derfor pralede alle mine nye kolleger meget af deres løn.}

\begin{tipbox}{本课发音与吞音要点（Udtale \& Reduktioner）}
\begin{itemize}
  \item \textbf{传闻表达 \textit{skulle efter sigende} 的隐藏吞音}：\textit{sigende} 拼写里有 \textit{g} 和 \textit{d}，但实际标准读音是 \ipa{ˈsiːənə}（\textbf{\textit{g} 完全不发音，\textit{nd} 中的 \textit{d} 也不发音}！）；整串搭配在口语里常读作 \ipa{sgu ˈɛfdɐ ˈsiːənə}。
  \item \textbf{法则与价值观核心词}：\textit{Janteloven} \ipa{ˈjandəˌlʌʊ̯ˀən}（\textit{-lov-} 读合口双元音 \ipa{ʌʊ̯ˀ}）；\textit{samfund} \ipa{ˈsɑmˌfɔnˀ}（词尾 \textit{d} 不发音）；\textit{beskedenhed} \ipa{beˈsgeːˀðənˌheːˀð}（软 d）。
  \item \textbf{推测情态词的弱读与重音}：作主观推测时，情态动词常轻读或次重读，重音落在后面的动词或表语上：\textit{Det må \textbf{være}} \ipa{de mʌ ˈvɛːɐ}、\textit{Det kan \textbf{ikke}} \ipa{de kʰa ˈeɡə}（口语极快时 \textit{kan} 缩为 \ipa{kʰa}）；\textit{oprindeligt} \ipa{ʌbˈʁenˀəli}（\textit{d} 与词尾 \textit{gt} 均不发音）。
\end{itemize}
\end{tipbox}

\begin{transbox}
\small
\textbf{\color{NCDTeal}[Tekst B 中文对照]}\\
\textbf{初入丹麦团队}：刚加入我的丹麦 IT 团队时，我原以为必须大声推销自己的想法才能被听见，因为我以前读到过竞争在跨国公司里很重要。[1]（\textbf{但丹麦办公室其实并不总是那样运作的。}）在我的第一次项目会议上，我注意到老板说话非常平和，并且听取实习生意见的认真程度不亚于对待资深开发者。此外，大家在休息时一起喝咖啡、聊周末计划，毫不炫耀。这一定是《詹代法则》至今仍在影响团队协作的一个明显例子。尽管如今这条法则据说已经没那么明显，但谦逊低调在办公室里依然很有生命力。
\end{transbox}

% ---------------------------------------------------------------------
\ncdsubsec{三、Gloser（核心生词表）}

\begin{tipbox}{词表重点}
情态动词 \textit{at burde} 的现在时是特殊的 \textit{bør}（不规则无 \textit{-r} 以外的词干变音！），过去时是 \textit{burde}。另外注意 \textit{at skille sig ud}（脱颖而出）是反身动词短语，不要漏掉反身代词 \textit{sig / mig / dig}；\textit{grundlæggende} 作定语和表语时词形始终不变。
\end{tipbox}

\begin{glosertable}
动词 [har] & at prale (af) & \textit{(praler, pralede, har pralet)} & 炫耀，吹嘘（接介词 \textit{af}） \\
动词 [har] & at skille sig ud & \textit{(skiller, skilte, har skilt sig ud)} & 脱颖而出，显得与众不同 \\
动词 [har] & at opmuntre (til) & \textit{(opmuntrer, opmuntrede, har opmuntret)} & 鼓励，激励 \\
动词 [har] & at behandle & \textit{(behandler, behandlede, har behandlet)} & 对待，处理 \\
动词 [har] & at se ned på & \textit{(ser, så, har set ned på)} & 看不起，轻视 \\
动词 [har] & at ændre (sig) & \textit{(ændrer, ændrede, har ændret)} & 改变，变化 \\
情态动词 [har] & at burde & \textit{(bør, burde, har burdet)} & 应当，按理说 \\
\midrule
名词 & et samfund & \textit{(samfundet, samfund, samfundene)} & 社会 \\
名词 & en lighed & \textit{(ligheden, ligheder, lighederne)} & 平等；相似之处 \\
名词 & en værdi & \textit{(værdien, værdier, værdierne)} & 价值观；价值 \\
名词 & en beskedenhed & \textit{(beskedenheden, ingen flertal)} & 谦虚，低调 \\
名词 & en succes & \textit{(succesen, succeser, succeserne)} & 成功 \\
名词 & en iværksætter & \textit{(iværksætteren, iværksættere, iværksætterne)} & 创业者，企业家 \\
名词 & en bremse & \textit{(bremsen, bremser, bremserne)} & 阻碍，刹车 \\
名词 & en indflydelse & \textit{(indflydelsen, ingen flertal)} & 影响，影响力 \\
\midrule
形容词 & grundlæggende & \textit{(grundlæggende, 不变)} & 根本的，基础的 \\
副词短语 & efter sigende & \textit{(固定副词短语)} & 据说，传闻中 \\
\end{glosertable}

% ---------------------------------------------------------------------
\ncdsubsec{四、Grammatik \& Faste vendinger（核心语法与地道表达精讲）}

\begin{grammatikbox}{语法点 1：情态动词的“主观推测”与“传闻转述”功能}
在 A2 阶段我们学过情态动词的“能力/义务”本义（\textit{kan} 会、\textit{skal} 必须、\textit{må} 允许）。进入 B1 阶段，同一批情态动词用来表达说话人对事实真伪的\textbf{推断把握度}或\textbf{消息来源}（接不带 \textit{at} 的动词原形）：
\begin{itemize}
  \item \textbf{\da{må} + 原形（必然/高度确定，90\%+ 把握）}：根据现有线索合乎逻辑地推断“一定是……”。例：\textit{Hun har boet her i ti år, så hun \da{må kende} Janteloven.}（她在这儿住了十年，肯定了解詹代法则。）
  \item \textbf{\da{kan} + 原形（可能性推测，约 50\% 把握）}：表示“有可能会……”。例：\textit{Reglerne \da{kan virke} strenge for udlændinge.}
  \item \textbf{\dabold{kan ikke} + 原形（绝对不可能）}：断然否定某种可能性。例：\textit{Det \dabold{kan ikke passe}, at alle tænker ens!}
  \item \textbf{\dabold{skal / skulle (efter sigende)} + 原形（传闻转述“据说”）}：表明信息来自他人或舆论，\textbf{说话人自己不作担保}！用过去时形式 \dabold{skulle} 比用 \textit{skal} 更有“姑且一听、保持距离”的客观感。例：\textit{Byen \dabold{skulle efter sigende være} meget åben i dag.}
  \item \textbf{对“过去事件”的推测：\da{må / kan / skulle} + \dabold{have} + 过去分词}：表示“一定已经…… / 据说曾经……”。例：\textit{Kulturen \da{må have ændret sig} gennem de sidste år.}
\end{itemize}
\end{grammatikbox}

{\small
\renewcommand{\arraystretch}{1.25}
\begin{xltabular}{\linewidth}{@{} >{\bfseries\color{NCDNavy}\raggedright\arraybackslash}p{2.9cm} >{\raggedright\arraybackslash}p{3.2cm} >{\raggedright\arraybackslash}p{3.2cm} >{\raggedright\arraybackslash}X @{}}
\toprule
\textbf{情态结构} & \textbf{推测/语义色彩} & \textbf{把握度 / 立场} & \textbf{B1 例句} \\
\midrule
\endfirsthead
\toprule
\textbf{情态结构} & \textbf{推测/语义色彩} & \textbf{把握度 / 立场} & \textbf{B1 例句} \\
\midrule
\endhead
\bottomrule
\endfoot
\textit{må} + 原形 & 逻辑必然推测\newline “一定，肯定” & 高度肯定（约 90\%+） & \textit{Det \da{må være} svært at begynde i et nyt firma.} \\
\midrule
\textit{kan} + 原形 & 开放性可能\newline “可能，也许” & 半信半疑（约 50\%） & \textit{Det \da{kan skabe} en god balance på kontoret.} \\
\midrule
\textit{kan ikke} + 原形 & 逻辑排除\newline “不可能” & 坚决排除 & \textit{Det \dabold{kan ikke være} hele forklaringen.} \\
\midrule
\textit{bør / burde} + 原形 & 合理推断或建议\newline “理应，按说应该” & 期待/规范判断 & \textit{Man \da{bør} også rose unge iværksættere.} \\
\midrule
\textit{skulle (efter sigende)} + 原形 & 传闻转述\newline “据说，听说” & 客观转述，不担保真假 & \textit{Den nye chef \dabold{skulle efter sigende være} meget direkte.} \\
\end{xltabular}
}

\begin{grammatikbox}{语法点 2：中英文对比排雷——三大推测致命误区}
\begin{enumerate}
  \item \textbf{“一定不是……”绝不能用 *\textit{må ikke}！}：很多同学把“一定是 \textit{må være}”直接加 \textit{ikke} 变成 *\textit{Det må ikke være sandt.}——错！在丹麦语中，\dabold{\textit{må ikke}} 永远只表示\textbf{“禁止/不许”}（等同于英语 \textit{must not}）。表达推测“一定不是 / 不可能是”，必须彻底换成 \dabold{\textit{kan ikke være}}！
  \item \textbf{传闻“据说”绝不逐字翻译成 *\textit{det er sagt, at...}！}：英语有 \textit{is said to be}，但丹麦语最地道的对应表达就是情态动词 \dabold{\textit{skal / skulle (efter sigende) være}}（也可以用 \textit{man siger, at...}，但 \textit{skulle efter sigende} 是标准 B1/B2 书面高分标志）。
  \item \textbf{情态动词后面绝不加 \textit{at}！}：哪怕中间插入了 \textit{efter sigende}，后面的不定式依然是光杆原形：\textit{Han skulle efter sigende \dabold{bo}} i Aarhus（绝不是 *\textit{at bo}！）。
\end{enumerate}
\end{grammatikbox}

% ---------------------------------------------------------------------
\ncdsubsec{五、Mønsterøvelser（句型操练）}

\textbf{练习 A：语境选词——用 \textit{må, kan, kan ikke, bør} 或 \textit{skulle} 填空（注意推测与传闻语义）}
\begin{enumerate}
  \item Når han arbejder 60 timer om ugen, \underline{\hspace{1.8cm}} han være meget træt om aftenen.
  \item Det \underline{\hspace{2.2cm}} passe, at hun allerede har glemt aftalen --- hun ringede jo for ti minutter siden! (否定)
  \item Den nye restaurant i midtbyen \underline{\hspace{2.0cm}} efter sigende have den bedste fisk i byen.
  \item Janteloven \underline{\hspace{1.8cm}} både have positive og negative sider for unge medarbejdere.
  \item Vi \underline{\hspace{1.8cm}} give nye kolleger mere plads til at fortælle om deres idéer.
  \item Hun svarer ikke på telefonen; hun \underline{\hspace{1.8cm}} allerede være taget hjem fra arbejde.
  \item Kulturen på arbejdspladser \underline{\hspace{2.0cm}} efter sigende have ændret sig meget de sidste ti år.
  \item Du \underline{\hspace{1.8cm}} ikke prale for meget på dit første jobmøde, hvis du vil skabe et godt indtryk.
\end{enumerate}

\textbf{练习 B：语体改写——将下列直接陈述改写为带情态推测或传闻转述的 B1 句子}
\begin{enumerate}
  \item 用 \textit{må være} 改写：\textit{Jeg er helt sikker på, at det er svært at starte i et nyt land.} $\rightarrow$ \textit{Det...}
  \item 用 \textit{skulle efter sigende være} 改写：\textit{Folk siger, at kulturen i København er mere international i dag.} $\rightarrow$ \textit{Kulturen i København...}
  \item 用 \textit{kan ikke passe} 改写：\textit{Det er helt umuligt, at Janteloven er forsvundet fuldstændigt.} $\rightarrow$ \textit{Det...}
  \item 用 \textit{må have + 分词} 改写（过去推测）：\textit{Jeg er sikker på, at hun har glemt sin paraply på kontoret.} $\rightarrow$ \textit{Hun...}
  \item 用 \textit{kan skabe} 改写：\textit{Måske skaber for meget konkurrence stress på arbejdspladsen.} $\rightarrow$ \textit{For meget konkurrence...}
  \item 纠错改写：\textit{Det må ikke være rigtigt, at han tjener så mange penge i et studiejob!} $\rightarrow$
\end{enumerate}

\textbf{练习 C：中译丹（精准使用 \textit{må, kan (ikke), skulle efter sigende}）}
\begin{enumerate}
  \item 初来乍到时，这一定很难习惯，因为大家都很低调。
  \item 据说这所大学在创新领域非常开放。
  \item 他不可能已经回家了，因为他的包还放在桌子上呢。
  \item 一点点谦虚可以营造出一个更加平等的合作氛围。
  \item 他们昨天一定在办公室里聊得很愉快。
\end{enumerate}

% ---------------------------------------------------------------------
\ncdsubsec{六、DU3 Modul 3 实战微任务与参考答案}

\begin{emailbox}{Skrivning Opgave 2 实战微任务（要求 100 词以上）}
\textbf{Situation}: Din ven Daniel overvejer at søge et job i en dansk virksomhed, men han har hørt om ”Janteloven” og er bange for, at han ikke må vise sine kompetencer.\\
\textbf{Opgave}: Skriv en mail til Daniel. Forklar ham, hvordan Janteloven fungerer på en moderne arbejdsplads, hvad han bør passe på, og hvordan han alligevel kan vise sin erfaring. Brug mindst tre forskellige modale udtryk (f.eks. \textit{må være, kan være, skulle efter sigende, bør}). Skriv \textbf{mindst 100 ord}.
\end{emailbox}

\begin{answerbox}
\textbf{Før du læser 参考回答}
\begin{enumerate}
  \item \textit{På den ene side kan Janteloven skabe et trygt arbejdsmiljø, hvor chefen og medarbejderne behandler hinanden ligeværdigt. På den anden side kan den være en bremse, hvis unge bliver bange for at skille sig ud.}
  \item \textit{Skribenten blev overrasket over, at chefen talte meget roligt og lyttede lige så meget til praktikanterne som til de erfarne udviklere. Det skyldtes blandt andet, at Janteloven stadig præger samarbejdet på kontoret.}
\end{enumerate}

\textbf{Tekst B 段落逻辑补全答案（Modul 3 Læsning Opgave 3 风格）}
\begin{enumerate}
  \item[] \da{B}（前文说“我原以为必须大声推销自己的想法”，下文突然转折描写丹麦老板很平静、倾听所有人意见，因此必须用转折句 \textit{Men sådan fungerer et dansk kontor altså ikke altid.} 承上启下；\textit{A} 语义完全违背事实，\textit{C} 和 \textit{D} 与上下文逻辑脱节）
\end{enumerate}

\textbf{练习 A 答案}
\begin{enumerate}
  \item \da{må}（每周工作 60 小时，逻辑必然“一定很累”）
  \item \da{kan ikke}（十分钟前才打过电话，“不可能”忘记）
  \item \da{skulle / skal}（\textit{skulle efter sigende} 据说/传闻，用过去时 \textit{skulle} 更显客观存疑）
  \item \da{kan}（表示客观可能性“可能有正反两面”）
  \item \da{bør}（提出合理规范建议“应当”）
  \item \da{må}（不接电话，高度推测“一定是已经回家了”）
  \item \da{skulle / skal}（与 \textit{efter sigende} 搭配表示传闻）
  \item \da{bør / må}（劝告/提醒“不应该/不许”）
\end{enumerate}

\textbf{练习 B 答案}
\begin{enumerate}
  \item \textit{Det \da{må være} svært at starte i et nyt land.}
  \item \textit{Kulturen i København \da{skulle efter sigende være} mere international i dag.}
  \item \textit{Det \da{kan ikke passe}, at Janteloven er forsvundet fuldstændigt.}
  \item \textit{Hun \da{må have glemt} sin paraply på kontoret.}
  \item \textit{For meget konkurrence \da{kan skabe} stress på arbejdspladsen.}
  \item \textit{Det \dabold{kan ikke være} rigtigt, at han tjener så mange penge i et studiejob!}（否定推测禁区：必须把 \textit{må ikke} 改为 \textit{kan ikke}）
\end{enumerate}

\textbf{练习 C 答案}
\begin{enumerate}
  \item \textit{Når man er ny, \da{må det være} svært at vænne sig til det, fordi alle er så beskedne.}
  \item \textit{Dette universitet \da{skulle efter sigende være} meget åbent for innovation.}
  \item \textit{Han \da{kan ikke} allerede være taget hjem, for hans taske står stadig på bordet.}
  \item \textit{En lille smule beskedenhed \da{kan skabe} et mere ligeværdigt samarbejde.}
  \item \textit{De \da{må have haft} en rigtig god samtale på kontoret i går.}
\end{enumerate}

\textbf{Modul 3 实战参考范文（123 词）}\\
\textit{Til: Daniel \quad|\quad Emne: Janteloven og dit nye job i Danmark}

\textit{Hej Daniel}

\textit{Jeg forstår godt din bekymring! Når man læser om Janteloven online, \da{må det næsten lyde}, som om man slet ikke må være ambitiøs i Danmark. Men det \dabold{kan ikke passe} i praksis på en moderne arbejdsplads.}

\textit{Janteloven \dabold{skulle efter sigende være} blevet meget svagere i dag, især i internationale firmaer. Her \da{kan} det tværtimod \da{være} et problem, hvis man er for stille omkring sine idéer. Det vigtige er bare måden, man gør det på: Du \da{bør} fokusere på, hvordan din viden kan hjælpe hele teamet i stedet for kun at prale af dig selv. Danske chefer sætter nemlig stor pris på ærlighed og godt samarbejde.}

\textit{Held og lykke med ansøgningen --- det skal nok gå rigtig godt!}

\textit{Mange hilsner\\
\mbox{[Dit navn]}}

\small\textbf{点评}：全篇正文 123 词，全面回应丹尼尔的顾虑并给出实用建议；恰当运用 \textit{må næsten lyde}、\textit{kan ikke passe}、\textit{skulle efter sigende være}、\textit{kan være} 与 \textit{bør} 五种情态用法，语气自然克制、逻辑清晰。
\end{answerbox}
